<%= heading2 %>{Abbreviations}

With \verb+\gls{...}+, you can enter abbreviations: the first time you call it, the long form is used.
When reusing \verb+\gls{..}+ the short form is automatically displayed.
The abbreviation is also automatically inserted in the abbreviation list.
With \verb+\glspl{...}+ the plural form is used.
If you want the short form to appear directly at the first use, you can use \verb+\glsunset{...}+ to mark an abbreviation as already used.
The opposite is achieved with \verb+\glsreset{...}+.
Both only change the ``already used'' flag: the effect shows at the next \verb+\gls{...}+ of that abbreviation, so they have to come before it in the text.
\verb+\glsresetall+ (e.g., after the abstract) lets every abbreviation appear in its long form again.

Abbreviations are defined in \verb+abbreviations.tex+ by means of \verb+\newabbreviation{key}{short}{long}+ (or \verb+\newacronym+ for an acronym that is read as a word).

More information at: \url{https://ctan.org/pkg/glossaries-extra}.

<%- bexample %>
At the first pass, the \gls{er} was 5.
At the second pass was \gls{er} 3.
The plural form can be seen here: \glspl{er}.
To demonstrate what the list of abbreviations looks like for longer description texts, \glspl{rdbms} must be mentioned here.

\gls{dante} is a local \TeX\ user group.
The German-speaking local \TeX\ user group is \gls{dante}.
A \gls{gp} is a medical doctor.
I went to my surgery to see the \gls{gp}.
<%- eexample %>

\todo{Include difference between acronym and abbreviation - \url{https://tex.stackexchange.com/a/154566/9075}}
